\documentclass[11pt]{article}
\usepackage[a3paper,landscape,margin=12mm]{geometry}
\usepackage{amsmath,amssymb,tikz,graphicx,adjustbox,booktabs,caption,hyperref}
\usetikzlibrary{calc,backgrounds}
\definecolor{highlight}{RGB}{0,0,0}
\input{drawings/preamble.tex}
\input{drawings/legend.tex}
\pdfinfoomitdate=1
\pdftrailerid{}
\pdfsuppressptexinfo=15
\hypersetup{pdftitle={Pauli semiwebs of the double-checking circuits},
 pdfauthor={},pdfsubject={},pdfkeywords={},pdfcreator={},pdfproducer={},
 colorlinks=true,allcolors=highlight,bookmarksopen=false}
\setlength{\parindent}{0pt}
\setlength{\parskip}{4pt}
\begin{document}
\begin{semiwebaddition}
\pdfbookmark[0]{Reading the diagrams}{guide}
{\LARGE Pauli semiwebs of the double-checking circuits}\par
\medskip
{\large A basis for the $d=3$ and $d=5$ ZX diagrams}
\bigskip

\begin{minipage}[t]{.47\linewidth}
\textbf{Contents}

Each following page shows one basis element. Products of the 90 elements
at $d=3$, or the 264 elements at $d=5$, generate every semiweb of the
corresponding graph, up to overall Pauli phase. The basis is not unique.
The counts depend on the chosen ZX graph.

\medskip
\begin{tabular}{lrr}
\toprule
 & Count & Pages\\
\midrule
$d=3$: incoming $X_q,Z_q$ &14&2--15\\
$d=3$: additional output support &14&16--29\\
$d=3$: no boundary support &62&30--91\\
$d=5$: incoming $X_q,Z_q$ &38&92--129\\
$d=5$: additional output support &38&130--167\\
$d=5$: no boundary support &188&168--355\\
\bottomrule
\end{tabular}

\medskip
``No boundary support'' means that the labels on all input and output
wires are $I$. These semiwebs may still have defects. The full Pauli-web
subspaces have dimensions 12 and 37; their subspaces with no boundary
support have dimensions 6 and 19.

\medskip
\textbf{Circuit}

The $d=5$ headings F1--F7 and U1--U7 identify the original fold and
unfold CNOT groups. Gates within a group use separate drawing columns
to avoid overlapping connections.

\medskip
\textbf{Sources}

The definitions and drawing styles follow Kissinger and van de Wetering,
\emph{ZX-Flow}, \href{https://arxiv.org/abs/2603.09580}{arXiv:2603.09580},
Sec.~3. The double-checking circuits are from the construction of Gidney,
Shutty and Jones, \href{https://arxiv.org/abs/2409.17595}{arXiv:2409.17595}.
The $d=5$ graph retains the mixed $T/T^\dagger$ pattern of the SOFT source.
\end{minipage}\hfill
\begin{minipage}[t]{.49\linewidth}
\textbf{Reading a diagram} (diagram from \url{https://github.com/tqec/tqec})

\begin{center}
\raisebox{-.5\height}{\semiwebinput[.40\linewidth]{drawings/examples/teleportation_left.tikz}}%
\hfill${}={}$\hfill
\raisebox{-.5\height}{\semiwebinput[.055\linewidth]{drawings/examples/teleportation_T.tikz}}%
\hfill${}={}$\hfill
\raisebox{-.5\height}{\semiwebinput[.40\linewidth]{drawings/examples/teleportation_right.tikz}}
\end{center}
Both circuits give $T$, up to a scalar, with an $S^\dagger$ correction on
the right. The overlays show different semiwebs.

\textbf{Legend}\quad
\semiwebinput[.84\linewidth]{drawings/examples/reading_legend.tikz}

Up to normalisation:
\mbox{\semiwebCaptionGate{\frac{\pi}{4}} ${}=T$ ($-\pi/4$: $T^\dagger$)};
\mbox{\semiwebCaptionGate{\frac{\pi}{2}} ${}=S$ ($-\pi/2$: $S^\dagger$)};
\mbox{\semiwebCaptionCnot ${}={\rm CNOT}$};
measurement outcome $0$: \mbox{\semiwebCaptionEffect{Z} ${}=\langle+|$ ($X$ basis)},
\mbox{\semiwebCaptionEffect{X} ${}=\langle0|$ ($Z$ basis)};
preparations \mbox{\semiwebCaptionState{X} ${}=|0\rangle$}
and \mbox{\semiwebCaptionState{Z} ${}=|+\rangle$};
resources \mbox{\semiwebCaptionPhaseState{\frac{\pi}{4}} ${}=|T\rangle$},
\mbox{\semiwebCaptionPhaseState{\frac{\pi}{2}} ${}=|Y\rangle$}.

At every spider, the incident Pauli labels satisfy
\[
 w_v|v(\theta_v)\rangle
 =\lambda_v|v(\theta_v+\delta_v)\rangle.
\]
A nonzero phase shift $\delta_v$ modulo $2\pi$ is a \emph{defect}.
No defects gives a Pauli web. Numbered violet stars mark the affected
nodes; the phase shifts are listed below each diagram. The phases
printed inside the spiders are the original phases.
A $\pi$ defect at a preparation or measurement changes
$|+\rangle$ or $\langle+|$ to $|-\rangle$ or $\langle-|$.

The boundary labels and scalar satisfy
\[
 K P_{\rm in}=\lambda P_{\rm out}K_{\boldsymbol\delta}.
\]
$K$ represents the selected double-checking stage, and $K_{\boldsymbol\delta}$
has the marked phases changed. Multiplication includes the Pauli scalar
phases, and phase changes must be composed. Adding two defect shifts
computed on the original spider is not generally correct.

Each TikZ file records the affected node, qubit and phase change in
comments. The star is positioned relative to that named node.

\end{minipage}
\vfill
\small The drawings suppress common ZX normalisation factors, which cancel between the two sides of each identity.
\newpage
\input{drawings/pages.tex}
\end{semiwebaddition}
\end{document}
